\scanpage{167}
\mspage{4}
% archive sheet number: 166
% A vertical emphasis line stands at the left of the lemma and its parenthesis.
\uline{Lemme} 1.2.6. Pour que l'opération \inserted{(1.2.4)} du groupe d'inertie $I$ sur $\pi_1^{\unsure{(p')|(p)}}(X_{\bar K},\xi)\deleted{^{\unsure{(p)}}}$ soit ``\uline{modérément} ramifiée'' i.e. pour que l'image \inserted{de $I$ dans $\operatorname{Aut}\pi_1^{(p')}(X_{\bar K},\xi)$} soit un gpe profini \uline{premier à $p$}, il \unsure{f. et s.} qu'il en soit ainsi de l'image de $I$ dans $\operatorname{Aut\,ext}(\pi_1^{(p')}(X_{\bar K},\xi))$ (via (1.1.5)). En particulier, cette condition ne dépend pas du choix du point \deleted{\unsure{géom. $\xi$}} choisi de $X_K(K)$.

\deleted{\unsure{Si cette condition est}}
\deleted{On dit alors, par abus de}
(Si cette dernière condition est satisfaite, on dit aussi que $I$ \uline{opère modérément} (sous-entendu : mod. aut. intérieurs) sur $\pi_1^{(p')}(X_{\bar K},\xi)$.)

Ceci posé, nous avons en vue la démonstration des deux résultats principaux :

% slanted
\margin{\deleted{\illegible{5} \unsure{d'un cas}}}
\uline{Théorème} 1.3. \inserted{(Théorème de l'action essentiellement modérée du groupe de monodromie \inserted{locale})}
Soit $X$ de type fini sur \deleted{\unsure{le corps}} $K$, $\xi$ un pt géom de $X$, $\bar K$ la clôture séparable correspondante de $K$, d'où $I=\operatorname{Gal}(\bar K/K)$ opérant \inserted{``extérieurement''} sur $\pi_1^{(p')}(X_{\bar K},\xi)$.
\deleted{\unsure{il existe un sous-groupe}}
\inserted{Alors (pour mémoire) $\pi_1^{(p')}(X_{\bar K},\xi)$ est} topologiquement \deleted{à engendrement} \inserted{de type} fini, et \deleted{\unsure{il}} il existe un sous-groupe ouvert $I'$ de $I$ qui opère \uline{modérément} (1.2.6) sur $\pi_1^{(p')}(X_{\bar K},\xi)$.

\uline{Remarque} 1.3.1. En fait, utilisant le résultat précédent, \inserted{(ou au choix le résultat 1.4 dans le cas de la dim. relative 1)} \deleted{\unsure{la théorie relative}} \inserted{et la résolution des singularités dans le cas non propre,} un argument dû à J. P. Murre (SGA 1 \uline{XIII}) permet de prouver ceci : Soit $X$ de t. f. sur $k$ sép clos, \deleted{$\xi$ un pt de \unsure{car.}} d'exp. car. $p$, $\xi$ un pt géom
